\chapter{Lưu trữ: Volume, Claim và trạng thái trên một node}
\label{ch:storage}

\section{Chồng trừu tượng}

Hệ thống file của container chết cùng nó. Kubernetes tách \emph{yêu
cầu} lưu trữ khỏi việc \emph{cung cấp} nó qua ba object:

\begin{description}
  \item[PersistentVolume (PV)] một mẩu lưu trữ thật --- một thư mục,
  một volume EBS, một export NFS.
  \item[PersistentVolumeClaim (PVC)] yêu cầu của workload: ``5\,Gi,
  ReadWriteOnce.'' Claim là thứ pod mount; nó gắn (bind) vào một PV.
  \item[StorageClass] nhà máy: cách tạo PV theo nhu cầu khi một claim
  xuất hiện (\emph{cấp phát động}).
\end{description}

Sự gián tiếp này là tính khả chuyển: \code{postgres-course.yaml} chỉ
nói ``tôi cần 5\,Gi''; còn \emph{làm thế nào} nó thành đĩa là việc của
cụm. Cùng manifest đó nhận một thư mục cục bộ trên k3s và một volume
EBS trên EKS.

\section{local-path trên k3s}

k3s đi kèm StorageClass \code{local-path}: một PVC trở thành một thư mục
dưới \code{/var/lib/rancher/k3s/storage/} trên SSD của node. Đơn giản,
nhanh, và thẳng thắn về giới hạn --- dữ liệu nằm trên \emph{đĩa của node
đó}. Trên cụm nhiều node, điều đó ghim pod vào node; trên cụm một node
này, nó chỉ có nghĩa đúng như nó nói.

\begin{hardware}
local-path cung cấp \emph{bền vững} (persistence: dữ liệu sống sót qua
khởi động lại pod và node), không phải \emph{bền chắc} (durability: một
SSD hỏng là mất tất cả) và không phải \emph{sao lưu} (một lệnh
\code{DELETE} sai lầm không có bản sao nào để quay về).
Mảnh còn thiếu, nói thẳng, là một CronJob chạy \code{pg\_dump} lên MinIO
hoặc ra ngoài máy. Quy tắc ngón tay cái đáng ghi nhớ: cơ sở dữ liệu
không có khôi phục đã kiểm chứng chỉ là một cái cache.
\end{hardware}

\section{Hai kiểu claim trong manifest}

MinIO (một Deployment với một replica) dùng PVC độc lập; các cơ sở dữ
liệu dùng \code{volumeClaimTemplates} bên trong StatefulSet --- một PVC
\emph{đúc riêng cho từng replica} và, quan trọng nhất, \emph{gắn lại vào
pod thay thế} chứ không phải vào một pod sống thứ hai. Sự phân biệt đó
là nửa lưu trữ của cạm bẫy hỏng dữ liệu ở Chương~15.

\begin{lstlisting}[language=yaml]
volumeClaimTemplates:
  - metadata: { name: data }
    spec:
      accessModes: [ReadWriteOnce]     (*@\co{1}@*)
      resources: { requests: { storage: 5Gi } }
# in the pod template:
volumeMounts:
  - { name: data, mountPath: /var/lib/postgresql/data }
env:
  - { name: PGDATA, value: /var/lib/postgresql/data/pgdata }  (*@\co{2}@*)
\end{lstlisting}

\begin{description}
  \item[\co{1}] \code{ReadWriteOnce} = mount đọc-ghi bởi một
  \emph{node} tại một thời điểm --- vừa khít cho cơ sở dữ liệu.
  (\code{RWX} --- nhiều node --- cần lưu trữ kiểu NFS và là chế độ
  người ta muốn cho upload dùng chung, trước khi phát hiện lưu trữ đối
  tượng, Chương~9, là câu trả lời tốt hơn.)
  \item[\co{2}] Mẹo PGDATA, một kinh điển thực thụ: image Postgres từ
  chối khởi tạo vào thư mục không rỗng, mà volume ext4 sinh ra đã có
  \code{lost+found}. Trỏ PGDATA sâu thêm một thư mục là né được. Bạn sẽ
  gặp đúng mẫu này trong hàng năm trời StackOverflow về
  Postgres-trên-Kubernetes.
\end{description}

\begin{pitfall}[xóa StatefulSet, dữ liệu vẫn còn --- hoặc mất]
Hai bất ngờ trong một. Thứ nhất: xóa StatefulSet \emph{không} xóa PVC
của nó --- theo thiết kế, dữ liệu sống lâu hơn sai lầm điều phối; tạo
lại StatefulSet là nó gắn lại. Thứ hai, bẫy ngược: xóa PVC khi đang thử
nghiệm \emph{có} xóa PV và thư mục của nó (reclaimPolicy
\code{Delete}, mặc định của cấp phát động). Sự bất đối xứng này đáng
thuộc lòng trước khi bạn dọn một namespace: object workload là thứ dùng
xong bỏ; claim mới là nơi dữ liệu nằm.
\end{pitfall}

\begin{handson}[nhìn xuyên qua lớp trừu tượng]
\begin{lstlisting}[language=bash]
kubectl -n ts get pvc                       # claims and their PVs
kubectl get pv                              # cluster-scoped, note the paths
ssh ts-server 'sudo ls /var/lib/rancher/k3s/storage/'
# pvc-...  one directory per PV -- your databases, as plain files
\end{lstlisting}
Giải ảo chính là mục đích: một PVC trên cụm này là một thư mục. Trên
EKS, cùng hai lệnh kubectl đó sẽ hiện ID volume EBS --- cùng object,
khác nhà máy.
\end{handson}

\section{Ngoài phạm vi dự án}

Trên Kubernetes có quản lý, StorageClass ánh xạ sang đĩa cloud
(EBS/gp3 trên EKS, Persistent Disk trên GKE) với snapshot, mã hóa và
đổi cỡ --- object \code{VolumeSnapshot} đưa sao lưu vào cùng mô hình
khai báo. Nhưng mẫu công nghiệp mà chương này hướng tới lại là phép
trừ: các đội ngày càng giữ \emph{cơ sở dữ liệu ngoài cụm} (RDS, Cloud
SQL) và để Kubernetes chỉ chạy những thứ phi trạng thái --- đổi sự chăm
chút StatefulSet ở Chương~15 lấy một hóa đơn dịch vụ có quản lý. Giờ
bạn đã thấy cả hai mặt của cuộc đánh đổi đó.

\begin{quiz}
\begin{enumerate}
  \item PV, PVC, StorageClass: pod mount cái nào, cái nào thuộc phạm vi
  cụm, và sự gián tiếp mua được gì cho manifest của dự án trên một EKS
  tương lai?
  \item Vì sao \code{ReadWriteOnce} là đủ cho Postgres, và loại
  workload nào đẩy người ta về phía RWX --- thường là sai?
  \item Giải thích cả hai chiều của sự bất đối xứng khi xóa: StatefulSet
  mất/PVC còn, PVC mất/dữ liệu mất.
  \item Bền vững, bền chắc và sao lưu mỗi thứ nghĩa là gì, và local-path
  cung cấp thứ nào?
\end{enumerate}
\end{quiz}
